% Transcripción de las páginas preliminares del impreso de 1972 (páginas PDF 1-6), conservada como
% registro; no forma parte del lector (FRONT_MATTER_ES.tex tiene la portada y los créditos de esta
% edición).

% SGA 7, tomo I — traducción española del frontispicio y la introducción.
% UTF-8 sin BOM; finales de línea LF. Fragmento; no se afirma una compilación autónoma.
% Fuente francesa corriente: source_layers_r01_20260820/corrected_candidate/
%   SGA7I_FRONT_MATTER_FR_CORRECTED_CANDIDATE.tex
% Testigo rector: páginas PDF 1--12 / índices de escaneo 0--11.
% Los avisos de laguna, obstrucción y procedencia son editoriales, no texto histórico.

\providecommand{\SGASevenSourceGapES}[1]{\textit{[laguna visual: #1]}}
\providecommand{\SGASevenSourceObstructionES}[1]{\textit{[oculto por #1]}}

% SGA7I-FRONT-ES-P000 — página PDF rectora 1 — cubierta.
\begin{center}
Lecture Notes in Mathematics

Colección de informes informales y seminarios

Editada por A. Dold, Heidelberg, y B. Eckmann, Zürich

\bigskip
288

\bigskip
Seminario de Geometría Algebraica de Bois-Marie

1967--1969

SGA 7 I

Dirigido por A. Grothendieck

con la colaboración de M. Raynaud y D. S. Rim

\bigskip
Grupos de monodromía

en geometría algebraica

\bigskip
Springer-Verlag

Berlin : Heidelberg : New York
\end{center}
\clearpage

% SGA7I-FRONT-ES-P001 — página PDF rectora 2 — guarda editorial.
% El ex libris de Nottingham oculta el centro de la presentación de la editorial.
% Los fragmentos ingleses visibles se conservan sin completar: traducirlos exigiría inventar el texto oculto.
\noindent
\textit{Fragmentos visibles de la presentación editorial:}\par
This series aims to report new d\SGASevenSourceObstructionES{el ex libris de la biblioteca}s - quickly, informally\par
and at a high level. The type of m\SGASevenSourceObstructionES{el ex libris de la biblioteca}\par
1. Preliminary drafts of original p\SGASevenSourceObstructionES{el ex libris de la biblioteca}\par
2. Lectures on a new field, or pre\SGASevenSourceObstructionES{el ex libris de la biblioteca}\par
3. Elaboraciones de seminario\par
4. Informes de reuniones, siempre que \SGASevenSourceObstructionES{el ex libris de la biblioteca}\par
\hspace*{1em}a) sean de interés excepcional o \SGASevenSourceObstructionES{el ex libris de la biblioteca}\par
\hspace*{1em}b) estén dedicadas a un solo tema.\par
Texts which are out of print but s\SGASevenSourceObstructionES{el ex libris de la biblioteca}thin these categories.\par
The timeliness of a manuscript i\SGASevenSourceObstructionES{el ex libris de la biblioteca}finished or tentative.\par
Thus, in some instances, proofs \SGASevenSourceObstructionES{el ex libris de la biblioteca}ch have been or will\par
se publicarán posteriormente en otro lugar.\par
Los manuscritos deberán comprender no \SGASevenSourceObstructionES{el ex libris de la biblioteca}\par
La publicación de \textit{Lecture Notes} tiene por objeto \SGASevenSourceObstructionES{el ex libris de la biblioteca}la comunidad científica, en\par
el hecho de que una editorial comercial, Spri\SGASevenSourceObstructionES{el ex libris de la biblioteca}documentos que de otro modo\par
tendrían una difusión restringida \SGASevenSourceObstructionES{el ex libris de la biblioteca}documentados en la\par
literatura científica. -\par

\begin{center}
Autor: LECTURES\par
Signatura: S/QA36.3\par
Núm. de libro: 539502\par

UNIVERSITY OF NOTTINGHAM

George Green Library

of Science and Engineering
\end{center}

\noindent
Manuscritos\par
Los manuscritos se reproducen mediante un procedimiento fotográfico; por tanto, deben mecanografiarse con sumo cuidado.\par
Los símbolos que no figuren en la máquina de escribir deberán añadirse a mano con tinta negra indeleble. Las correcciones al texto\par
mecanografiado deberán hacerse pegando el texto enmendado sobre el antiguo o cubriendo los errores con líquido corrector blanco.\par
Los autores reciben 75 ejemplares gratuitos.\par
El texto mecanografiado se reduce ligeramente durante la reproducción; no se obtendrán los mejores resultados si el texto de una\par
página rebasa el límite total de 18 x 26,5 cm (7 x 10\ensuremath{\tfrac12} pulgadas). A petición de los interesados, la editorial facilitará\par
papel especial en el que esté delimitada el área de mecanografiado.\par
Los manuscritos en inglés, alemán o francés deberán enviarse al Prof. Dr. A. Dold, Mathematisches Institut der Universität Heidelberg,\par
69 Heidelberg/Alemania, Tiergartenstraße, o al Prof. Dr. B. Eckmann, Eidgenössische Technische Hochschule,\par
CH-8006 Zürich/Suiza.\par

\medskip
\noindent
Lecture Notes in Physics\par
Publicados hasta la fecha\par
% Los títulos y datos de catálogo se conservan literalmente como datos bibliográficos históricos.
Vol. 1: J. C. Erdmann, Wärmeleitung in Kristallen, theoretische Grundlagen und fortgeschrittene experimentelle Methoden. II, 283 Seiten. 1969. DM 20,--\par
Vol. 2: K. Hepp, Théorie de la renormalisation. III, 215 pages. 1969. DM 18,--\par
Vol. 3: A. Martin, Scattering Theory: Unitarity, Analyticity and Crossing. IV, 125 pages. 1969. DM 14,--\par
Vol. 4: G. Ludwig, Deutung des Begriffs »physikalische Theorie« und axiomatische Grundlegung der Hilbertraumstruktur der Quantenmechanik durch Hauptsätze des Messens. XI, 469 Seiten. 1970. DM 28,--\par
Vol. 5: M. Schaaf, The Reduction of the Product of Two Irreducible Unitary Representations of the Proper Orthochronous Quantummechanical Poincaré Group. IV, 120 pages. 1970. DM 14,--\par
Vol. 6: Group Representations in Mathematics and Physics. Edited by V. Bargmann. V, 340 pages. 1970. DM 24,--\par
Vol. 7: R. Balescu, J. L. Lebowitz, I. Prigogine, P. Résibois, Z. W. Salsburg, Lectures in Statistical Physics. V, 181 pages. 1971. DM 18,--\par
Vol. 8: Proceedings of the Second International Conference on Numerical Methods in Fluid Dynamics. Edited by M. Holt. IX, 462 pages. 1971. DM 28,--\par
Vol. 9: D. W. Robinson, The Thermodynamic Pressure in Quantum Statistical Mechanics. V, 115 pages. 1971. DM 14,--\par
Vol. 10: J. M. Stewart, Non-Equilibrium Relativistic Kinetic Theory. III, 113 pages. 1971. DM 14,--\par
Vol. 11: O. Steinmann, Perturbation Expansions in Axiomatic Field Theory. III, 126 pages. 1971. DM 14,--\par
\clearpage

% SGA7I-FRONT-ES-P002 — página PDF rectora 3 — portada con fichas de préstamo.
% Solo se traducen o conservan los fragmentos visibles; no se reconstruye el título oculto.
\begin{center}
Lecture Notes in Mathematics

Colección de informes informales y seminarios

Editada por A. Dold, Heidelberg, y B. Eckmann, Zürich

\bigskip
288

\bigskip
Groupe\SGASevenSourceObstructionES{las fichas de préstamo de la biblioteca}

en Géor\SGASevenSourceObstructionES{las fichas de préstamo de la biblioteca}

\bigskip
Springer-Verlag

Berlin : Heidelberg : New York 1972
\end{center}
\clearpage

% SGA7I-FRONT-ES-P003 — página PDF rectora 4 — página de copyright y datos editoriales.
% NOTA EDITORIAL, FUERA DEL AVISO HISTÓRICO:
\begin{quote}
\textit{Nota editorial: el aviso de copyright reproducido a continuación se conserva únicamente como procedencia histórica; no restringe esta edición matemática independiente ni su publicación. El copyright no constituye un requisito de permiso, seguridad, revisión jurídica, autorización de derechos, traducción, auditoría, compilación ni publicación. No se formula ninguna declaración general de dominio público ni de licencia abierta.}
\end{quote}

\noindent
Clasificaciones temáticas AMS (1970): Primaria: 14D05, 14D15, 14K99\par
\hspace*{8em}Secundaria: 13D10, 14E20, 14F20\par

\medskip
\noindent
ISBN 3-540-05987-3 Springer-Verlag Berlin - Heidelberg - New York\par
ISBN 0-387-05987-3 Springer-Verlag New York - Heidelberg - Berlin\par

\medskip
\noindent
Esta obra está sujeta a copyright. Quedan reservados todos los derechos, tanto sobre la totalidad como sobre una parte del material de que se trat\SGASevenSourceGapES{borde derecho cortado}\par
en particular los de traducción, reimpresión, reutilización de ilustraciones, radiodifusión, reproducción mediante máquinas fotocopiadora\SGASevenSourceGapES{borde derecho cortado}\par
o medios similares, y almacenamiento en bancos de datos.\par
Conforme al § 54 de la Ley alemana de Copyright, cuando se hagan copias para un uso distinto del privado deberá abonarse una tasa a la editoria\SGASevenSourceGapES{borde derecho cortado}\par
cuya cuantía se determinará mediante acuerdo con la editorial.\par

\medskip
\noindent
© by Springer-Verlag Berlin : Heidelberg 1972. Library of Congress Catalog Card Number 72-88730. Printed in Germany.\par

\medskip
\noindent
Offsetdruck: Julius Beltz, Hemsbach/Bergstr.\par
\clearpage

% SGA7I-FRONT-ES-P004 — página PDF rectora 5 — título interno del seminario.
\begin{center}
SEMINARIO DE GEOMETRÍA ALGEBRAICA DE BOIS-MARIE

1967--1969

\bigskip
GRUPOS DE MONODROMÍA EN GEOMETRÍA ALGEBRAICA

(SGA 7 I)

\bigskip
Dirigido por A. GROTHENDIECK

con la colaboración de

M. RAYNAUD y D. S. RIM
\end{center}
\clearpage

% SGA7I-FRONT-ES-P005 — página PDF rectora 6 — procedencia de Internet Archive.
\begin{center}
Digitalizado por Internet Archive

en 2022 con financiación de

Kahle/Austin Foundation

\bigskip
\texttt{https://archive.org/details/groupesdemonodro0288unse}
\end{center}
\clearpage

